% Einheit 5 von 5 – Besondere Linien im Dreieck und Satz des Thales (bank/winkel-dreiecke/e5.jsonl, zusammenbau v0.1)
%% TODO Zweigzeile Teil 1: was hier gelernt wird (Satz in Schülersprache aus dem Einheitstitel) – Einheit 5
\einheitenkopf[e5]{Einheit 5 von 5 · Besondere Linien im Dreieck und Satz des Thales}
\zweigzeile{ab Kl. 6, je nach Buch bis Kl. 7 · P10}
\setcounter{aufgabe}{22}

%% TODO Titel: Ich-kann-Satz fehlt in der Bank (Platzhalter: Kettenname) – Nr. 23
%% TODO Anweisung: ein Satz über der ersten Teilaufgabe fehlt in der Bank – Nr. 23
\begin{aufgabe}{Besondere Linien}
\begin{teile}
\teil Eine Linie geht durch die Ecke C und steht senkrecht auf der Seite AB. Welche Linie im Dreieck ABC ist das? \\ \kreuz{Mittelsenkrechte} \\ \kreuz{Winkelhalbierende} \\ \kreuz{Höhe} \\ \kreuz{Seitenhalbierende}
\end{teile}
\swa{Verbinde A und B und konstruiere mit dem Zirkel die Mittelsenkrechte der Strecke AB. In welchem Punkt schneidet sie AB? M( \leerfeld | \leerfeld )}{\begin{ksys}[xmin=0,xmax=10,ymin=0,ymax=6] \punkt{2}{2}{A} \punkt{8}{4}{B} \end{ksys}}
\swa{Verbinde A und B und konstruiere mit dem Zirkel die Mittelsenkrechte der Strecke AB. In welchem Punkt schneidet sie AB? M( \leerfeld | \leerfeld )}{\begin{ksys}[xmin=0,xmax=10,ymin=0,ymax=6] \punkt{1}{1}{A} \punkt{7}{5}{B} \end{ksys}}
\swa{Verbinde A und B und konstruiere mit dem Zirkel die Mittelsenkrechte der Strecke AB. In welchem Punkt schneidet sie AB? M( \leerfeld | \leerfeld )}{\begin{ksys}[xmin=0,xmax=10,ymin=0,ymax=6] \punkt{3}{5}{A} \punkt{9}{1}{B} \end{ksys}}
\swa{Verbinde A und B und konstruiere mit dem Zirkel die Mittelsenkrechte der Strecke AB. In welchem Punkt schneidet sie AB? M( \leerfeld | \leerfeld )}{\begin{ksys}[xmin=0,xmax=10,ymin=0,ymax=6] \punkt{1}{4}{A} \punkt{7}{0}{B} \end{ksys}}
\swa{Verbinde A und B und konstruiere mit dem Zirkel die Mittelsenkrechte der Strecke AB. In welchem Punkt schneidet sie AB? M( \leerfeld | \leerfeld )}{\begin{ksys}[xmin=0,xmax=10,ymin=0,ymax=6] \punkt{2}{5}{A} \punkt{6}{1}{B} \end{ksys}}
\swfrage{Was bleibt gleich, was ändert sich?}
\swa{Zeichne das spitzwinklige Dreieck ABC. Konstruiere zwei Mittelsenkrechten und den Umkreis. Gib den Mittelpunkt U an. U( \leerfeld | \leerfeld )}{\begin{ksys}[xmin=0,xmax=10,ymin=0,ymax=8] \punkt{1}{1}{A} \punkt{7}{1}{B} \punkt{3}{5}{C} \end{ksys}}
\swa{Konstruiere mit dem Zirkel die Winkelhalbierende von α. Wie groß ist jeder der beiden Teilwinkel? \leerfeld °}{\winkel[5]{64}{\alpha}}
\swa{Zeichne das Dreieck ABC. Konstruiere zwei Winkelhalbierende und den Inkreis. Gib den Mittelpunkt I an. I( \leerfeld | \leerfeld )}{\begin{ksys}[xmin=0,xmax=10,ymin=0,ymax=10] \punkt{1}{1}{A} \punkt{7}{1}{B} \punkt{1}{9}{C} \end{ksys}}
\swa{Zeichne das Dreieck ABC und mit dem Geodreieck die Höhe von C auf die Seite AB, wenn nötig auf ihre Verlängerung. Gib den Fußpunkt D an. D( \leerfeld | \leerfeld )}{\begin{ksys}[xmin=0,xmax=10,ymin=0,ymax=8] \punkt{1}{1}{A} \punkt{7}{1}{B} \punkt{4}{5}{C} \end{ksys}}
\swa{Zeichne das Dreieck ABC und seine drei Seitenhalbierenden. Gib den Schwerpunkt S an. S( \leerfeld | \leerfeld )}{\begin{ksys}[xmin=0,xmax=10,ymin=0,ymax=8] \punkt{1}{1}{A} \punkt{7}{1}{B} \punkt{4}{7}{C} \end{ksys}}
\swa{Zeichne eine Strecke AB von 8 cm Länge und darüber den Thaleskreis. Wähle einen Punkt C darauf und zeichne das Dreieck ABC. Welchen Radius hat der Halbkreis, und wo liegt der rechte Winkel? r = \leerfeld[cm]}{\winkelstrahl[10][A]}
%% TODO Darstellung für schwach fehlt in der Bank (winkel-dreiecke-e5-k1-s8-v1; Abschnitt „Für schwache Schüler“)
\swz[3]{Die Punkte A und B liegen sich auf einem Kreis mit dem Mittelpunkt M gegenüber, die Strecke AB geht durch M. C ist ein weiterer Punkt auf dem Kreis. Begründe, dass das Dreieck ABC bei C einen rechten Winkel hat.}{}
\swz[3]{Im Viereck ABCD schneiden sich die Diagonalen AC und BD im Punkt M. Es gilt MA = MC = MD = 4 cm. Begründe, dass der Winkel des Vierecks bei D ein rechter ist.}{}
\end{aufgabe}

%% TODO Titel: Ich-kann-Satz fehlt in der Bank (Platzhalter: Kettenname) – Nr. 24
%% TODO Anweisung: ein Satz über der ersten Teilaufgabe fehlt in der Bank – Nr. 24
\begin{aufgabe}{Besondere Linien – Fehler finden, Begründen, Anwendung}
\swz[3]{Lukas soll im Dreieck ABC die Höhe von C auf AB zeichnen. Er verbindet C mit der Mitte von AB. Finde den Fehler.}{}
\swfrage{Warum ist der Schnittpunkt der Mittelsenkrechten eines Dreiecks ABC von allen drei Ecken gleich weit entfernt? Begründe.}
\swa{Drei Dörfer A, B und C wollen einen gemeinsamen Funkmast, der von allen drei Dörfern gleich weit entfernt ist (1 Kästchen = 1 km). Konstruiere den Standort. Gib seine Koordinaten und die Entfernung zu den Dörfern an. ( \leerfeld | \leerfeld )}{\begin{ksys}[xmin=0,xmax=10,ymin=0,ymax=8] \punkt{1}{1}{A} \punkt{9}{1}{B} \punkt{1}{7}{C} \end{ksys}}
\end{aufgabe}

\uebersichtskasten{Mittelsenkrechte: steht senkrecht auf der Mitte einer Seite; alle drei schneiden sich im Umkreismittelpunkt – er ist von allen Ecken gleich weit entfernt. \\ Winkelhalbierende: teilt den Winkel in zwei gleiche Hälften; alle drei schneiden sich im Inkreismittelpunkt – er ist von allen Seiten gleich weit entfernt. \\ Höhe: Lot von einer Ecke auf die gegenüberliegende Seite (beim stumpfwinkligen Dreieck auf ihre Verlängerung). Seitenhalbierende: von der Ecke zur Seitenmitte; ihr Schnittpunkt ist der Schwerpunkt, er teilt jede im Verhältnis 2 : 1. \\ Thales: Liegt C auf dem Halbkreis über der Strecke AB (AB ist der Durchmesser), ist der Winkel bei C ein rechter – und umgekehrt. \\ Rechtwinkliges Dreieck mit Hypotenuse 6 cm: Halbkreis mit Radius 3 cm um die Mitte von AB, C darauf wählen.}
